\chapter{Горизонтальная реконструкция вакуумного гессиана}
\label{ch:v8-gauge-invariant-vacuum-hessian-reconstruction-gate}

Провал голдстоуновского критерия не требует выбрасывать всю локальную
квадратичную информацию. Следовая метрика $K_{\Phi}$ канонически задаёт
ортогональное разложение transfer-касательного пространства на нарушенную
калибровочную орбиту и горизонтальный срез. Для независимого орбитального
репера $X$ точный грамов блок особенно прост:
\begin{equation}
 G_{\rm orb}=X^{\mathsf T}K_{\Phi}X=14I_3,
 \qquad
 P_{\rm orb}=XG_{\rm orb}^{-1}X^{\mathsf T}K_{\Phi},
 \qquad Q=I_{30}-P_{\rm orb}.
 \label{eq:v8-horizontal-orbit-projector}
\end{equation}
Он удовлетворяет
\begin{equation}
 P_{\rm orb}^2=P_{\rm orb},\qquad
 P_{\rm orb}^{\mathsf T}K_{\Phi}=K_{\Phi}P_{\rm orb},\qquad
 \rank Q=27.
 \label{eq:v8-horizontal-projector-properties}
\end{equation}

Канонический квадратичный quotient прежнего фиксированно-фонового блока
определим формулой
\begin{equation}
 H_{\rm quot}=Q^{\mathsf T}H_{\Phi}Q,
 \qquad H_{\rm quot}X_{\rm orb}=0.
 \label{eq:v8-horizontal-quotient-hessian}
\end{equation}
Теперь голдстоуновский критерий выполнен точно. При этом
\begin{equation}
 \rank H_{\rm quot}=26,\qquad
 \dim\ker H_{\rm quot}=4=3_{\rm Goldstone}+1_{\rm horizontal}.
 \label{eq:v8-horizontal-kernel-decomposition}
\end{equation}
Последняя мода не является остатком калибровочной орбиты. В координатах
полного вещественного transfer-кадра её можно выбрать как
$v_0=(0,6,0,0,0,0,0,-1,0,1,0,\ldots,0)^{\mathsf T}$; она
$K_{\Phi}$-ортогональна орбите и зануляется уже исходным грамовым
гессианом.

После той же кинетической нормировки точные моменты равны
\begin{equation}
 \Tr(M_{\Phi,{\rm quot}}^4)=\frac{1118917}{882},
 \qquad
 N_{\rm bos}^{\rm BV}
 =\Tr(M_{\Phi,{\rm quot}}^4)+3\Tr(M_B^4)
 =\frac{226371884}{159201}.
 \label{eq:v8-horizontal-bv-bosonic-ledger}
\end{equation}
С уже выведенным хиральным фермионным вкладом получается
\begin{equation}
 N_{\rm quad}^{\rm BV}
 =\frac{226371884}{159201}-92
 =\frac{211725392}{159201}.
 \label{eq:v8-horizontal-full-quadratic-numerator}
\end{equation}

\begin{theorem}[Канонический квадратичный BV-quotient]
Следовая метрика единственно задаёт ортогональный проектор на нарушенную
калибровочную орбиту. Горизонтальная компрессия прежнего гессиана
зануляет ровно три голдстоуновских направления и оставляет одну
некалибровочную плоскую моду.
\end{theorem}

\begin{nogocriterion}[Квадратичная реконструкция не равна полному родителю]
Формула $Q^{\mathsf T}H_{\Phi}Q$ задаёт корректную вторую вариацию на
выбранном следовой метрикой горизонтальном срезе, но сама по себе не строит
гладкий нелинейный калибровочно-инвариантный функционал. Поэтому
$211725392/159201$ является точным квадратичным BV-кандидатом, а не
безусловным физическим коэффициентом.
\end{nogocriterion}

\begin{protocol}
\begin{enumerate}
 \item орбитальная метрика: \textbf{$14I_3$};
 \item горизонтальная размерность: \textbf{$27$};
 \item ранг quotient-гессиана: \textbf{$26$};
 \item ядро: \textbf{три голдстоуна и одна горизонтальная мода};
 \item квадратичный BV-числитель: \textbf{$211725392/159201$};
 \item нелинейный калибровочно-инвариантный родитель: \textbf{не выведен};
 \item следующий гейт: \textbf{родительский подъём горизонтальной плоской
 моды}.
\end{enumerate}
\end{protocol}

Машинный сертификат сохранён в
\path{s2t_v8_gauge_invariant_vacuum_hessian_reconstruction_gate_results.json}.